\documentclass{article}
\usepackage[T2A]{fontenc}
\usepackage[utf8]{inputenc}
\usepackage[russian]{babel}
\usepackage{amsmath,amssymb}
\usepackage{graphicx}
\usepackage{hyperref}
\title{Синтез биологических и архитектурных форм\\ методом GRA-обнуления (Bio-Arch GRA)}
\author{О. Б. Битсоев\thanks{Независимый исследователь, \texttt{https://orcid.org/0009-0004-1872-1153}}}
\date{\today}

\begin{document}
\maketitle

\begin{abstract}
Представлен эксперимент по биоинспирированному архитектурному дизайну на основе многоуровневой системы GRA-обнуления (Geometric Recursive Analytics). Параметрическая модель колонны с внутренней пористостью оптимизируется так, чтобы одновременно удовлетворять биологическим критериям (лёгкость, высокая удельная прочность) и архитектурным требованиям (устойчивость, технологичность, эстетические пропорции). Оптимизация минимизирует сумму функционалов «пены», вычисляемых на основе простых механических расчётов и пропорциональных правил. Результаты показывают, что GRA порождает структурно эффективные и визуально гармоничные формы, напоминающие трабекулярную костную ткань в колонне.
\end{abstract}

\section{Введение}
GRA-обнуление \cite{gra_metatheorem} — методология устранения артефактов предметных областей путём рекурсивного подавления междисциплинарных расхождений (пены). В предыдущей работе метод был применён к текстовому синтезу (квантовая поэзия \cite{gra_quantum_poetry}). Здесь мы распространяем подход на геометрический синтез биологических и архитектурных форм.

\section{Метод}
\subsection{Уровни GRA и пена}
Определим два уровня:
\begin{align}
\Phi_{\text{bio}} &= w_1 \left(1 - \frac{F_{\text{cr}}}{F_{\text{opt}}}\right)^{\!2}
                  + w_2 \left(\frac{m}{m_{\text{opt}}}-1\right)^{\!2}
                  + w_3 \left(\frac{p}{p_{\text{opt}}}-1\right)^{\!2}, \\
\Phi_{\text{arch}} &= v_1 \left(1 - \frac{F_{\text{cr}}}{F_{\text{req}}}\right)^{\!2}
                   + v_2 \max\!\left(0,\frac{t_{\text{min}}-t}{t_{\text{min}}}\right)^{\!2}
                   + v_3 \left(\frac{H}{2R} - \phi\right)^{\!2},
\end{align}
где $F_{\text{cr}}$ — эйлерова критическая сила, $m$ — масса, $p$ — пористость, $t$ — минимальная толщина стенки, $\phi = 1.618$. Общая пена:
$J = \lambda \Phi_{\text{bio}} + (1-\lambda)\Phi_{\text{arch}}$.

\subsection{Механическое моделирование}
Колонна — цилиндр радиуса $R$ и высоты $H$ со сферическими порами. Критическая нагрузка определяется как
$F_{\text{cr}} = \pi^2 E I_{\text{min}} / H^2$, где $I_{\text{min}}$ — наименьший осевой момент инерции сечения, вычисляемый аналитически по координатам пор.

\subsection{Оптимизация}
Используется генетический алгоритм (библиотека DEAP). Геном кодирует $R$, $H$ и до 10 пор (признак активности, координаты, радиус). Размер популяции 50, число поколений 100, стандартные операторы скрещивания и мутации.

\section{Результаты}
После оптимизации наилучшая форма (рис.~\ref{fig:column}) достигла значений $\Phi_{\text{bio}}=0.12$, $\Phi_{\text{arch}}=0.09$ при $\lambda = 0.5$. Получена пористость $0.47$, $F_{\text{cr}}$ в пределах 90\% от целевого значения, отношение $H/D \approx 1.62$.

\begin{figure}[h]
\centering
\includegraphics[width=0.6\textwidth]{bio_column.png}
\caption{Оптимизированная био-архитектурная колонна с внутренними порами.}
\label{fig:column}
\end{figure}

\section{Заключение}
GRA-обнуление успешно объединило биологическую и архитектурную оптимизацию, породив лёгкую, устойчивую и эстетичную колонну. В дальнейшем параметрический генератор будет заменён вариационным автоэнкодером сеток для получения более богатых топологий.

\bibliographystyle{plain}
\bibliography{references}
\end{document}